\documentclass[11pt]{article}
\usepackage[T1]{fontenc}
\usepackage[utf8]{inputenc}
\usepackage[a4paper,margin=25mm]{geometry}
\usepackage{amsmath,amssymb,mathtools}
\usepackage{bm}
\usepackage{booktabs}
\usepackage{longtable}
\usepackage{array}
\usepackage{xcolor}
\usepackage{hyperref}
\usepackage{listings}

\hypersetup{colorlinks=true,linkcolor=blue!50!black,urlcolor=blue!60!black}

\newcommand{\vect}[1]{\bm{#1}}
\newcommand{\dd}{\,\mathrm{d}}
\newcommand{\code}[1]{\texttt{\detokenize{#1}}}
\newcommand{\rank}{\mathrm{rank}}
\newcommand{\owned}{\mathcal O}
\newcommand{\ghost}{\mathcal G}
\newcommand{\faces}{\mathcal F}
\newcommand{\eps}{\varepsilon}
\newcommand{\perm}{\pi}

\lstset{
  basicstyle=\ttfamily\small,
  breaklines=true,
  frame=single,
  columns=fullflexible
}

\title{DiscoGMPI: Distributed-Memory Implementation of a Three-Dimensional DG Maxwell Solver}
\author{DiscoGMPI project}
\date{August 2026}

\begin{document}
\maketitle

\begin{abstract}
This report documents the distributed-memory implementation currently present
in DiscoGMPI.  The focus is not to describe an idealized design, but the
solver architecture as implemented in the repository: tetrahedral
discontinuous Galerkin (DG) Maxwell operators, owned-plus-halo mesh storage,
MPI face-trace exchange, periodic distributed coupling, distributed
diagnostics, restart/checkpoint files, and production drivers for cavity,
periodic, nonlinear-PML, and scattering/RCS workflows.  The intended audience
is computational mathematicians and computational physicists who need to
understand how the mathematical discretization is mapped onto rank-local data
structures and MPI communication.
\end{abstract}

\tableofcontents

\section{Scope and Code Basis}

The implementation described here is the code currently wired through
\path{src/DiscoGMPI.jl}.  The public module includes the serial DG kernels,
distributed mesh construction, distributed DG wrappers, heterogeneous
material support, nonlinear PML support, distributed diagnostics, ParaView
output, checkpoint/restart support, and run-layout utilities.  The principal
source files inspected for this report are
\begin{itemize}
  \item \path{src/DistributedMesh3D.jl},
        \path{src/DistributedDG.jl},
        \path{src/DistributedIO.jl},
        and \path{src/TraceMaps.jl};
  \item \path{src/solver/data_containers.jl} and the kernels under
        \path{src/solver/kernels/};
  \item \path{src/PhysicalCoverage.jl}, \path{src/NonlinearPML.jl},
        \path{src/MaxwellDiagnostics.jl}, \path{src/MaxwellOutput.jl},
        \path{src/RunLayout.jl}, and \path{src/MaxwellDriverSupport.jl};
  \item the production and validation drivers under \path{examples/};
  \item MPI and validation tests under \path{test/};
  \item mesh-generation and post-processing scripts under
        \path{examples/meshes/}, \path{examples/}, and
        \path{benchmark/}.
\end{itemize}
There is no separate \path{drivers/} directory in the inspected tree; the
production entry points are Julia scripts in \path{examples/}.  The file
\path{src/rhs.jl} contains an older lower-level MPI RHS prototype, but it is
not included by \path{src/DiscoGMPI.jl}; therefore it is not part of the
current module load path documented as the production solver.

\section{High-Level Architecture}

DiscoGMPI is organized around a deliberate layering:
\begin{enumerate}
  \item A serial tetrahedral DG kernel constructs topology, geometry,
        reference operators, trace maps, flux faces, and Maxwell RHS kernels.
  \item A distributed mesh layer selects rank-local owned elements and a
        one-face ghost halo from a global partition.
  \item A distributed DG wrapper converts the local owned-plus-halo mesh into
        an ordinary \code{DGDiscretization}, then adds persistent MPI trace
        buffers and neighbor face metadata.
  \item Time-integration drivers call the same Maxwell RHS kernels on each
        rank after refreshing ghost face traces.
  \item Diagnostics and output integrate or write owned elements only, then
        use MPI reductions or rank-local files plus root-owned manifests.
\end{enumerate}
This architecture keeps most numerical kernels independent of MPI.  A rank
sees a local \code{RawVTUMesh} containing owned and ghost tetrahedra, and the
serial DG routines assemble local volume and face operators for that mesh.
Distributed correctness is enforced by ownership: ghost unknowns are used to
provide neighbor traces but ghost residuals and ghost diagnostic
contributions are removed.

\begin{longtable}{p{0.29\textwidth}p{0.64\textwidth}}
\caption{Primary implementation modules relevant to distributed Maxwell runs.}
\label{tab:module-map}\\
\toprule
\textbf{File} & \textbf{Implemented responsibility}\\
\midrule
\endfirsthead
\toprule
\textbf{File} & \textbf{Implemented responsibility}\\
\midrule
\endhead
\path{src/DistributedMesh3D.jl} &
Owned/ghost element selection, local node and element maps, boundary tags,
material IDs, face neighbor records, and MPI neighbor communication tables.\\
\path{src/DistributedDG.jl} &
\code{DistributedDGDiscretization}, local DG construction from distributed
meshes, persistent Maxwell send/receive buffers, ghost-trace exchange,
distributed RHS wrappers, distributed energy, and distributed time-step
wrappers.\\
\path{src/PhysicalCoverage.jl} &
Heterogeneous element material arrays, material-aware volume and surface
operators, material-aware distributed RHS functions, physical invariant
diagnostics, and distributed periodic Maxwell exchange.\\
\path{src/NonlinearPML.jl} &
Nodal nonlinear PML damping data, polynomial profiles, PML source terms, and
serial/distributed PML RHS and ESPRK/RK wrappers.\\
\path{src/DistributedIO.jl} &
Rank-local mesh-cache files, manifest/checksum validation, distributed
checkpoint/restart files, and run metadata/partition metadata.\\
\path{src/MaxwellDiagnostics.jl} &
Cubature-based distributed energy, field errors, component errors, charge,
linear momentum, angular momentum, optical chirality, and mesh-quality
metrics.\\
\path{src/MaxwellOutput.jl} &
Parallel VTK/PVTK field output, ParaView time-series collection files,
integration-point CSV output, and checkpoint pointer files.\\
\path{src/RunLayout.jl} &
Isolated run directory layout, configuration/input/status files, compatibility
links, and run-root naming utilities.\\
\path{examples/*.jl} &
Thin-to-medium production frontends for cavity, periodic, PML validation,
metallic-sphere scattering/RCS, NASA almond scattering/RCS, convergence, and
validation-matrix workflows.\\
\bottomrule
\end{longtable}

\section{Maxwell DG State and Local Discretization}

The electromagnetic state is stored as six nodal arrays,
\begin{equation}
  U_h = (E_x,E_y,E_z,H_x,H_y,H_z),
\end{equation}
where each component has dimensions $N_p\times K_{\mathrm{loc}}$ on a rank.
Here $N_p=(N+1)(N+2)(N+3)/6$ is the number of nodal basis functions on a
degree-$N$ tetrahedron, and $K_{\mathrm{loc}}$ includes owned and ghost
tetrahedra.  The concrete storage type is \code{MaxwellField}, defined in
\path{src/solver/data_containers.jl}; RHS storage is the analogous
\code{MaxwellRHS}.

The local DG discretization is built by \code{DGDiscretization} in
\path{src/solver/kernels/dg_discretization.jl}.  It contains:
\begin{itemize}
  \item the raw tetrahedral mesh \code{RawVTUMesh};
  \item element and boundary topology;
  \item geometric mappings and physical metrics;
  \item a \code{ReferenceTet} with nodal points, Vandermonde matrices,
        derivative matrices, weak derivative matrices, and mass matrix;
  \item reference face mass/lift operators;
  \item physical derivative operators on each affine tetrahedron;
  \item interior and boundary trace maps with face-node permutations;
  \item flux-face geometry and coloring information.
\end{itemize}
The serial kernels therefore retain the usual DG structure.  MPI is not
embedded in reference-element or physical-element operators; distributed
execution is obtained by presenting each rank with a consistent local mesh
and by exchanging just enough data to evaluate traces on partition interfaces.

\subsection{Semidiscrete Maxwell System}

On each element $K$, DiscoGMPI advances source-free Maxwell equations in the
form
\begin{equation}
  \eps_K \partial_t \vect E = \nabla\times\vect H,
  \qquad
  \mu_K \partial_t \vect H = -\nabla\times\vect E,
  \label{eq:maxwell}
\end{equation}
with positive scalar permittivity $\eps_K$ and permeability $\mu_K$.  The
homogeneous path passes scalar \code{epsilon} and \code{mu}; the
heterogeneous path uses \code{MaxwellElementMaterials} with one value per
rank-local element.

For affine tetrahedra, the element mass matrix is
\begin{equation}
  M_K = |J_K|M,
\end{equation}
where $M$ is the reference nodal mass matrix and $J_K$ is the affine mapping
Jacobian.  The mass-matrix energy used by the implementation is
\begin{equation}
  \mathcal H_h =
  \frac12\sum_K |J_K|\left[
    \eps_K\sum_{a\in\{x,y,z\}} E_{a,K}^T M E_{a,K}
    +
    \mu_K\sum_{a\in\{x,y,z\}} H_{a,K}^T M H_{a,K}
  \right].
  \label{eq:discrete-energy}
\end{equation}
In distributed runs the sum in \eqref{eq:discrete-energy} is performed over
owned elements only on each rank, followed by \code{MPI.Allreduce}.

\subsection{Hesthaven--Warburton and Poisson-Bracket Paths}

Two Maxwell formulations are represented by concrete types:
\code{HesthavenWarburtonFormulation} and \code{PoissonBracketFormulation}.
Both use the same field storage but different volume/surface operators.
The Hesthaven--Warburton path uses strong physical derivative operators for
the curl terms.  The Poisson-bracket path uses weak derivative matrices for
one partition and transposed weak derivative matrices for the complementary
partition, matching the implemented discrete skew pairing.

The Poisson-bracket surface operator supports centered and alternating flux
kinds.  Upwind flux is rejected for \code{PoissonBracketFormulation} because
the implemented upwind penalty is dissipative and does not define the
partitioned Poisson-bracket operator used by the ESPRK path.  The implemented
Poisson-bracket centered face correction is
\begin{equation}
  F_E^{PB}
  =-\frac{1}{2\eps}\vect n\times(\vect H^- - \vect H^+),
  \qquad
  F_H^{PB}
  =-\frac{1}{2\mu}\vect n\times(\vect E^- + \vect E^+).
  \label{eq:pb-centered-flux}
\end{equation}
For the alternating variant, the implementation selects
\begin{equation}
  \theta(\vect n)=
  \begin{cases}
  1, & n_x+n_y+n_z>0,\\
  0, & n_x+n_y+n_z<0,\\
  1/2, & n_x+n_y+n_z=0,
  \end{cases}
\end{equation}
sets $\alpha=1-\theta$, and uses
\begin{align}
  \widehat{\vect E} &= \theta \vect E^+ + \alpha \vect E^- ,\\
  F_E^{alt} &=
  -\frac{\alpha}{\eps}\vect n\times(\vect H^- - \vect H^+),\\
  F_H^{alt} &=
  -\frac{1}{\mu}\vect n\times\widehat{\vect E}.
\end{align}
Interior and distributed periodic alternating faces are implemented.  For
physical boundaries, the Poisson-bracket alternating code currently handles
PEC by returning no boundary contribution, and rejects PMC or absorbing
boundary kinds for this flux choice.

\section{Distributed Mesh Model}

\subsection{Global Partition Input}

The distributed mesh constructor accepts a global tetrahedral mesh and an
element-to-rank vector.  METIS-style element partitions are read by
\code{read\_metis\_epart}; internally ranks are zero-based to match MPI
rank IDs.  The constructor validates that the partition has exactly one entry
per tetrahedron and that all entries lie in the active rank range.

The implementation supports three practical setup modes:
\begin{enumerate}
  \item \code{build\_distributed\_dg}: every rank has the global mesh and
        partition available.
  \item \code{build\_distributed\_dg\_from\_root}: only a root rank provides
        the global mesh and partition; root builds rank-local mesh packets
        and sends them to the other ranks.
  \item \code{build\_distributed\_dg\_from\_partition}: each rank loads a
        previously written rank-local mesh shard from a distributed mesh
        cache.
\end{enumerate}
The production drivers expose this through \code{--distributed-mesh-dir},
\code{--rebuild-distributed-mesh}, and
\code{--collective-distributed-mesh-prep}.  The root-distribution path is
simple and robust but has a rank-zero setup bottleneck.  The collective mesh
preparation path trades memory during preprocessing for avoiding root-built
rank-local packets.  Once the distributed mesh cache exists, each rank reads
only its own serialized shard.

\subsection{Owned and Ghost Elements}

For a rank $r$, the owned global element set is
\begin{equation}
  \owned_r = \{K_i:\mathrm{owner}(K_i)=r\}.
\end{equation}
The ghost set is the one-face halo
\begin{equation}
  \ghost_r =
  \{K_j\notin\owned_r:\exists K_i\in\owned_r
  \text{ such that } K_i \text{ and } K_j
  \text{ share a triangular face}\}.
\end{equation}
This is constructed in \code{build\_ghost\_global\_elements\_face\_halo}.
Only a one-layer face halo is required because Maxwell surface fluxes depend
on immediate plus traces only.

The rank stores the sorted union $\owned_r\cup\ghost_r$, with separate local
index lists \code{partition.owned} and \code{partition.ghosts}.  This design
is important: ghost elements occupy ordinary element slots in the local
\code{DGDiscretization}, so the serial DG face-map machinery can identify
partition-interface faces as ordinary interior faces between an owned element
and a ghost element.  Ownership is then enforced outside the serial kernels.

\subsection{Local Node and Element Maps}

The local node set is the set of global vertices touched by all stored
owned and ghost tetrahedra.  The implementation stores
\begin{equation}
  \mathrm{local\_node}\leftrightarrow \mathrm{global\_node}
\end{equation}
maps in \code{LocalNodeGeometry}.  Element-level maps are stored in
\code{LocalElementTopology}: local vertices, global element IDs,
global-to-local element lookup, owner ranks, ownership flags, material IDs,
and per-face neighbor records.  The material ID vector is carried through
the distributed mesh so that heterogeneous material lookup remains local
after partitioning.

\subsection{Face Keys, Boundary Tags, and Non-Manifold Guards}

A tetrahedral face is keyed by sorting its three global vertex IDs.  A face
with one adjacent local record is treated as a boundary face, a face with two
records is an interior face, and any face with more than two records is
rejected as non-manifold.  Boundary condition tags are stored in a dictionary
from sorted global face key to integer boundary ID and are copied into
\code{FaceNeighborInfo.bc\_tag}.  This is used for PEC, PMC, absorbing, PML,
periodic, and scatterer-surface tags depending on the driver.

\section{Distributed DG Construction}

\code{build\_distributed\_dg\_from\_local\_mesh} converts a
\code{DistributedMesh} into a local \code{RawVTUMesh}.  The local mesh contains
all stored tetrahedra, including ghosts.  Boundary triangles are reconstructed
from faces that appear once in the rank-local stored mesh and are given the
corresponding boundary ID.  Faces between owned and ghost elements appear
twice, so they are not physical boundary faces in the local DG topology.

After the local \code{RawVTUMesh} is built, the ordinary serial
\code{DGDiscretization} constructor creates all reference, metric, trace, and
flux structures.  Finally, \code{build\_distributed\_maxwell\_exchange}
matches the distributed mesh communication table to the local DG face-node
arrays and allocates persistent send and receive buffers.

This construction means that most Maxwell kernels are blind to MPI.  A
distributed RHS call performs:
\begin{enumerate}
  \item face-trace exchange to refresh ghost face nodes;
  \item local Maxwell RHS evaluation using the serial DG kernels over the
        owned-plus-ghost local mesh;
  \item zeroing of all ghost RHS entries;
  \item owned-only use of the resulting RHS in time integration.
\end{enumerate}

\section{MPI Communication Strategy}

\subsection{Interface Metadata}

Partition-interface faces are represented by \code{DistributedDGInterfaceFace}.
Each entry stores the sorted global face key, neighbor rank, owned element and
face, owned face nodes, ghost element and face, and ghost face nodes.  Entries
are grouped by neighbor rank in \code{DistributedMaxwellExchange}.

For each neighbor $q$, the exchange object stores
\begin{equation}
  \code{send\_buffers}[q],\quad \code{recv\_buffers}[q]
  \in \mathbb R^{6N_{fp}N_{faces}(q)}.
\end{equation}
The factor six corresponds to $E_x,E_y,E_z,H_x,H_y,H_z$.  The buffers are
allocated once when the distributed DG object is built and reused on every
RHS call.  The send faces are sorted deterministically by global face key, so
both communicating ranks agree on buffer layout.

\subsection{Halo Exchange Algorithm}

The implemented Maxwell halo exchange is \code{exchange\_maxwell\_ghost\_traces!}.
Conceptually, for each neighbor rank $q$ it performs:
\begin{lstlisting}[language={},caption={Implemented Maxwell ghost-trace exchange.}]
for neighbor q
    pack owned face-node values of Ex,Ey,Ez,Hx,Hy,Hz into send_buffer[q]
end
for neighbor q
    post MPI.Irecv!(recv_buffer[q], source=q, tag=24017)
end
for neighbor q
    post MPI.Isend(send_buffer[q], dest=q, tag=24017)
end
MPI.Waitall(requests)
for neighbor q
    unpack recv_buffer[q] into ghost face-node slots
end
\end{lstlisting}
The implementation uses nonblocking receives and sends, but all requests are
waited on before RHS assembly starts.  Thus the current production path avoids
deadlock and has persistent buffers, but it does not yet overlap communication
with local volume or owned-interior surface work.

Only face-node values are required for partition-interface fluxes.  Ghost
element interiors are not synchronized.  This is consistent with the local
RHS flow because the physically relevant owned-element surface fluxes only
need plus traces on shared faces.  RHS values on ghost elements are filled by
the local serial call and then explicitly zeroed by
\code{zero\_ghost\_maxwell\_rhs!}.

\subsection{Distributed RHS Flow}

For homogeneous materials, the distributed RHS wrapper in
\path{src/DistributedDG.jl} is
\begin{equation}
  R_r(U_r) =
  \mathrm{zeroGhost}\left(
  R_{\mathrm{DG}}(\mathrm{exchange}(U_r))
  \right).
\end{equation}
For heterogeneous materials, \path{src/PhysicalCoverage.jl} defines analogous
overloads that pass \code{MaxwellElementMaterials} to material-aware volume,
interior-face, and boundary-face routines.  The material-aware distributed
RHS accepts optional \code{MaxwellBoundaryData}; this is used by scattering
drivers for time-dependent incident PEC data.

\subsection{Performance Consequences}

The communication volume per RHS evaluation scales like
\begin{equation}
  N_{\mathrm{send}} =
  6 N_{fp} \sum_{q\in\mathcal N_r} N_{\faces}(r,q),
\end{equation}
where $N_{\faces}(r,q)$ is the number of local owned faces adjacent to
neighbor $q$.  Since
$N_{fp}=(N+1)(N+2)/2$, surface communication grows quadratically with
polynomial order, whereas element storage grows cubically through $N_p$.
The one-face ghost halo adds memory overhead proportional to the partition
surface area; the per-rank metadata written by \code{write\_distributed\_run\_metadata}
records owned elements, ghost elements, ghost DOFs, neighbors, interface
faces, send/receive values, and halo ratios.

\section{Distributed Periodic Coupling}

Periodic problems need an additional communication layer because two physical
boundary faces may be paired even when they are not adjacent in the local
tetrahedral topology.  The implemented type is
\code{DistributedPeriodicMaxwellExchange} in \path{src/PhysicalCoverage.jl}.
The construction process is:
\begin{enumerate}
  \item each rank describes its owned boundary faces for the requested
        periodic boundary IDs;
  \item rank zero gathers descriptors and pairs minus/plus faces by centroid
        after applying the prescribed translation;
  \item face areas and face-node permutations are checked;
  \item the paired descriptors are broadcast to all ranks;
  \item each rank builds local \code{DistributedPeriodicFace} entries,
        including local partners when both faces are on the same rank and
        MPI-buffer offsets when the partner is remote.
\end{enumerate}

For a rectangular periodic box, the default boundary pairs are
\begin{equation}
  1\leftrightarrow 2,\qquad
  3\leftrightarrow 4,\qquad
  5\leftrightarrow 6,
\end{equation}
with translations $(-L_x,0,0)$, $(0,-L_y,0)$, and $(0,0,-L_z)$.  The
periodic mesh must be topologically compatible: translated partner faces
must have matching centroids, areas, and nodal locations within tolerance.
The earlier unstructured periodic mesh issue is consistent with this
requirement; if opposite surfaces are not translationally paired face by face,
the numerical periodic coupling is not the intended physical topology.

The periodic RHS first performs the ordinary ghost exchange, then performs
periodic trace exchange using tag \code{24018}.  Local periodic partners are
read directly from the local field arrays.  Remote partners are read from the
periodic receive buffers.  Periodic surface contributions are then added
before remaining physical boundary contributions are processed.

\section{Boundary Conditions and Physical Coverage}

Boundary conditions are imposed by constructing an exterior trace.  The
implemented boundary kinds are \code{MaxwellBC\_PEC},
\code{MaxwellBC\_PMC}, \code{MaxwellBC\_Absorbing}, and
\code{MaxwellBC\_None}.

For PEC boundaries, the exterior trace is
\begin{equation}
  \vect E^+ = -\vect E^- + 2(\vect n\cdot\vect E^-)\vect n,
  \qquad
  \vect H^+ = \vect H^-,
\end{equation}
so the averaged trace satisfies $\vect n\times\vect E=0$.  For PMC,
\begin{equation}
  \vect E^+ = \vect E^-,
  \qquad
  \vect H^+ = -\vect H^- + 2(\vect n\cdot\vect H^-)\vect n.
\end{equation}
The absorbing boundary trace zeroes the incoming characteristic using the
local impedance $Z=\sqrt{\mu/\eps}$ while retaining normal components.

The heterogeneous material path stores per-element $\eps_K$ and $\mu_K$ in
\code{MaxwellElementMaterials}.  For Hesthaven--Warburton fluxes, the
material-aware upwind flux uses minus and plus impedances.  For the
Poisson-bracket path, material-aware face fluxes call the implemented
Poisson-bracket surface flux with the minus-side material coefficients.

\section{Nonlinear PML Implementation}

The nonlinear PML implementation is contained in \path{src/NonlinearPML.jl}.
It stores nodal damping profiles
\begin{equation}
  \sigma_x,\quad \sigma_y,\quad \sigma_z
\end{equation}
as matrices matching the Maxwell field layout.  The helper
\code{polynomial\_pml\_sigma} evaluates one-sided polynomial damping profiles
that are zero in the interior and rise to \code{sigma\_max} at the outer
boundary.  \code{build\_six\_sided\_nonlinear\_pml} constructs Cartesian PML
slabs from distributed mesh bounds.

The PML appears as a nonlinear source added after the Maxwell spatial RHS.
At a node with electric field $\vect E$, magnetic field $\vect H$, and
damping vector $\vect\sigma$, the implementation computes auxiliary vectors
from $\vect E\times\vect H$ with denominators depending on the PML parameter
$a\in(0,1)$, $\eps$, $\mu$, and a positive regularization.  The resulting
electric and magnetic PML source components are added to the RHS.  Distributed
PML source addition is restricted to owned elements:
\begin{equation}
  R_r^{PML}(U_r) =
  \mathrm{zeroGhost}\left(
  R_r^{Maxwell}(U_r) + S_r^{PML}(U_r)\big|_{\owned_r}
  \right).
\end{equation}
Both homogeneous and heterogeneous-material PML RHS overloads are implemented,
as are distributed periodic PML variants.

For Poisson-bracket production paths, the code contains partitioned ESPRK PML
wrappers such as
\code{distributed\_maxwell\_nonlinear\_pml\_partitioned\_symplectic\_rk\_step!}.
Older helper routines for standard explicit RK PML time marching remain in
the code, but production Poisson-bracket drivers use ESPRK when PML is
enabled.

\section{Time Integration}

\subsection{Explicit RK}

\path{src/solver/kernels/time_integration.jl} implements classical explicit
RK methods of orders one through five, including a fifth-order
Dormand--Prince update without adaptive time stepping.  The workspace
\code{MaxwellRKWorkspace} stores the initial state, one stage state, and a
vector of stage RHS fields.

\subsection{Explicit Partitioned Symplectic RK}

The Poisson-bracket production drivers use explicit partitioned symplectic
Runge--Kutta (ESPRK) methods.  The scheme object stores
\begin{equation}
  a_i=\code{first\_weights}[i],
  \qquad
  b_i=\code{second\_weights}[i],
\end{equation}
and a first partition, either electric or magnetic.  The production Maxwell
drivers use magnetic-first stepping for Poisson-bracket Maxwell runs.  A
generic ESPRK substep computes the full RHS and applies it only to the active
partition:
\begin{lstlisting}[language={},caption={Implemented partitioned symplectic step.}]
for stage i
    if first_weight[i] != 0
        rhs!(rhs, U)
        U[first_partition] += dt * first_weight[i] * rhs[first_partition]
    end
    if second_weight[i] != 0
        rhs!(rhs, U)
        U[second_partition] += dt * second_weight[i] * rhs[second_partition]
    end
end
\end{lstlisting}
The timestep is estimated by \code{estimate\_maxwell\_dt} from element-size
diagnostics, polynomial order, material wave speed, and the user CFL value.
Production drivers then choose an integer number of steps so that the final
step lands exactly at \code{--final-time}.

\subsection{Nonautonomous Boundary Forcing}

The metallic-sphere and NASA-almond scattering drivers solve for scattered
fields with time-dependent incident PEC boundary data.  The incident boundary
data are stored in \code{MaxwellBoundaryData}, which contains a mutable
reference to the current physical time.  The ordinary ESPRK scheme object
does not store classical Runge--Kutta abscissas $c_i$.  For nonautonomous
forcing, the scattering driver therefore reconstructs the physical substep
clock from the cumulative composition weights.

For an H-first split, the driver keeps two clocks:
\begin{equation}
  t_H^{(0)}=t^n,\qquad t_E^{(0)}=t^n.
\end{equation}
Before a nonzero H-partition substep with weight $a_i$, it sets
\code{boundary\_data.time} to $t_H$, evaluates the RHS, updates the magnetic
partition, and then advances $t_H\leftarrow t_H+\Delta t\,a_i$.  Before a
nonzero E-partition substep with weight $b_i$, it sets the boundary time to
$t_E$, evaluates the RHS, updates the electric partition, and advances
$t_E\leftarrow t_E+\Delta t\,b_i$.  At the end of the step it sets the
boundary time to $t^{n+1}$.  This is implemented in
\code{timed\_partitioned\_symplectic\_rk\_step!} in
\path{examples/distributed_metallic_sphere_scattering.jl}.  The detail is
essential: evaluating incident fields only at integer step times would use
the wrong nonautonomous data inside the composition.

\section{Diagnostics and Invariants}

\path{src/MaxwellDiagnostics.jl} computes diagnostics by integrating over
owned elements only and reducing the resulting sums.  The quadrature
diagnostic uses Jaskowiec--Sukumar tetrahedral cubature, usually with order
\begin{equation}
  q=\max(2,2N+4)
\end{equation}
unless the user requests another cubature order.  The implemented diagnostic
quantities include:
\begin{itemize}
  \item electric, magnetic, and total energy;
  \item exact energy and energy error for analytical cases;
  \item $L^2$ norms and relative errors for $\vect E$, $\vect H$, and the
        combined field;
  \item per-component errors for $E_x,E_y,E_z,H_x,H_y,H_z$;
  \item energy-density norms and errors;
  \item optical chirality and optical-chirality density errors;
  \item electric and magnetic charge,
        $\int_\Omega \nabla\cdot(\eps\vect E_h)\dd x$ and
        $\int_\Omega \nabla\cdot(\mu\vect H_h)\dd x$;
  \item linear momentum based on $\eps\mu\,\vect E_h\times\vect H_h$;
  \item angular momentum based on
        $\vect x\times(\eps\mu\,\vect E_h\times\vect H_h)$;
  \item mesh-quality metrics for convergence CSVs.
\end{itemize}
The diagnostic values are not volume-normalized by default.  They are
physical integrals or $L^2$ norms over the computational domain, with the
normalization appearing only where an explicit relative error or RMS quantity
is computed.

The output schema for \code{diagnostics/quadrature\_diagnostics.csv} is
written by \code{write\_quadrature\_diagnostics\_header} and
\code{write\_quadrature\_diagnostics\_row}.  The energy history
\code{diagnostics/energy.csv} is written separately using the DG
mass-matrix energy.

\section{Output, Metadata, and Restart}

\subsection{Run Layout}

\path{src/RunLayout.jl} defines an isolated run directory layout:
\begin{equation}
\begin{array}{ll}
\code{config/} & resolved configuration, layout, run metadata, partition metadata,\\
\code{input/} & input manifest,\\
\code{diagnostics/} & CSV diagnostics,\\
\code{fields/} & ParaView PVTK/PVTU files,\\
\code{checkpoints/} & restart state shards,\\
\code{logs/},\code{plots/},\code{tables/} & downstream artifacts.
\end{array}
\end{equation}
The production drivers support either a fixed \code{--output-dir} or an
isolated timestamped directory under \code{--run-root} with a sanitized
\code{--run-name}.  Compatibility symlinks for common files such as
\code{energy.csv}, \code{quadrature\_diagnostics.csv}, and
\code{run\_metadata.toml} are created at the run root.

\subsection{Distributed Mesh Cache}

\code{write\_distributed\_mesh\_partition} writes one serialized mesh shard per
rank:
\begin{equation}
  \code{mesh\_rank0000.bin},\quad
  \code{mesh\_rank0001.bin},\quad \ldots
\end{equation}
Rank zero writes \code{mesh\_manifest.toml} only after all pieces are written.
The manifest records the format version, MPI rank count, global element
count, user metadata, per-rank file names, checksums, owned/ghost counts,
local nodes, and neighbor ranks.  \code{load\_distributed\_mesh\_partition}
checks the manifest, rank count, file existence, checksum, payload format,
payload version, and payload rank before returning the local mesh shard.

\subsection{Checkpoint and Restart}

\code{write\_distributed\_checkpoint} stores only owned Maxwell degrees of
freedom on each rank.  A checkpoint payload records the rank, MPI size, DG
order, nodes per element, owned global element IDs, and six field-component
arrays restricted to owned elements.  Rank zero writes
\code{checkpoint\_manifest.toml} with step, time, dt, metadata, file names, and
checksums.  Restart currently requires the same MPI rank count, DG order, and
element partition.  On load, the code verifies that the stored owned global
IDs match the current distributed mesh and then refreshes ghost traces.

\subsection{ParaView and Integration-Point Output}

Parallel field output uses \code{write\_parallel\_maxwell\_fields}.  Each rank
writes owned elements only as VTK Lagrange tetrahedra.  The fields include
numerical electric and magnetic fields, exact fields when supplied, field
errors, magnitudes, global element IDs, owner rank, polynomial order, and time.
The time series is indexed by \code{paraview\_series.csv} and
\code{fields.pvd}.  The \code{--paraview-every} driver option controls
snapshot cadence while still writing initial and final states in production
drivers.

Final integration-point output is written as one CSV per rank.  It contains
physical coordinates, quadrature weights, numerical and exact field values,
energy densities, chirality density, charge densities, linear momentum, and
angular momentum at cubature points of owned elements.

\section{Scattering and RCS as a Distributed Workflow}

Although this report focuses on the distributed Maxwell implementation, the
metallic-sphere and NASA-almond RCS drivers exercise several important
distributed features simultaneously: incident time-dependent PEC data,
scattered-field evolution, nonlinear PML, owned-surface integration, and MPI
reductions.

The metallic-sphere driver \path{examples/distributed_metallic_sphere_scattering.jl}
solves for the scattered field:
\begin{equation}
  \vect E_h=\vect E_{\mathrm{scat},h},\qquad
  \vect H_h=\vect H_{\mathrm{scat},h}.
\end{equation}
The incident plane wave is supplied by \path{src/MaxwellExactSolutions.jl}.
For a PEC scatterer surface, the total-field boundary condition is
\begin{equation}
  \vect n\times
  (\vect E_{\mathrm{scat}}+\vect E_{\mathrm{inc}})=0,
\end{equation}
which is imposed in the exterior trace as
\begin{equation}
  \vect n\times\vect E_{\mathrm{scat}}
  =-\vect n\times\vect E_{\mathrm{inc}}.
\end{equation}
The incident-aware exterior trace is implemented by
\code{pec\_incident\_boundary\_plus\_trace!} and routed through
\code{MaxwellBoundaryData}.

The RCS path samples owned PEC scatterer faces only.  It forms the total
magnetic field on the surface, computes the surface current
\begin{equation}
  \vect J_s = \vect n\times
  (\vect H_{\mathrm{scat},h}+\vect H_{\mathrm{inc}}),
\end{equation}
and accumulates a time-windowed Fourier sample at the incident angular
frequency.  For an observation direction $\widehat{\vect r}$, the implemented
spatial phase in the surface-current integral is
\begin{equation}
  \exp\left(-ik\,\widehat{\vect r}\cdot\vect x\right),
\end{equation}
and a Hann window is applied after \code{--rcs-start-time}.  The reduced
far-field amplitude is converted into
\begin{equation}
  \sigma(\theta,\phi)
  =
  4\pi
  \frac{|\vect E_\infty(\theta,\phi)|^2}{|\vect E_{\mathrm{inc}}|^2}.
\end{equation}
The output \code{diagnostics/rcs.csv} reports dimensional RCS values and dB
values for the requested angular pattern.  Sphere post-processing compares
against a Mie-series reference; almond post-processing compares against the
Austin CEM benchmark reference data when available.

\section{Convergence and Validation Infrastructure}

The convergence drivers are:
\begin{itemize}
  \item \path{examples/convergence_distributed_poisson_bracket_maxwell.jl}
        for PEC/PMC cavity eigenmodes;
  \item \path{examples/convergence_distributed_periodic_poisson_bracket_maxwell.jl}
        for periodic plane waves;
  \item \path{examples/convergence_distributed_periodic_isoresolution.jl}
        for isoresolution studies.
\end{itemize}
They write CSV convergence tables with aggregate and per-component errors,
rates, mesh-size information, and mesh-quality metrics.  Expected rates are
checked explicitly rather than merely checking monotone error reduction.

The validation-matrix orchestrator is
\path{examples/validate_distributed_maxwell_matrix.jl}.  It can run PEC,
PMC, periodic, and PML cases over selected rank counts and polynomial orders.
It launches MPI tests, compares rank-one outputs against multi-rank outputs
bitwise or within tolerance, checks convergence-rate policies, verifies
diagnostic schemas, checks conservative energy drift, checks PML energy
dissipation, and writes a machine-readable
\code{validation\_summary.json}.

The MPI regression tests include:
\begin{itemize}
  \item \path{test/test_mpi_metis_partition.jl}: partition reading,
        owned/ghost mesh construction, neighbor tables, deterministic face
        ordering, and lower-level trace exchange;
  \item \path{test/test_mpi_distributed_maxwell.jl}: distributed DG
        construction, mesh/cache/checkpoint I/O, ghost trace exchange,
        distributed RHS matching serial RHS, ghost RHS zeroing, energy
        matching serial energy, RK and ESPRK step matching, and PMC support;
  \item \path{test/test_mpi_face_permutation.jl}: face permutation
        correctness across interfaces;
  \item \path{test/test_mpi_physical_coverage.jl}: material and physical
        coverage paths;
  \item \path{test/test_mpi_nonlinear_pml.jl}: distributed nonlinear PML
        coverage.
\end{itemize}

\section{Memory and Scaling Considerations}

The main rank-local Maxwell storage is six dense matrices of size
$N_p\times K_{\mathrm{loc}}$ for the state, plus RHS and time-integration
workspace arrays.  For ESPRK, \code{MaxwellPartitionedRKWorkspace} stores a
single RHS field.  For classical explicit RK, \code{MaxwellRKWorkspace} stores
one initial field, one stage field, and one RHS field per stage.  Therefore
high-order classical RK has a larger memory footprint than ESPRK.

The distributed mesh duplicates ghost elements and their local nodes.  This
is a controlled overhead:
\begin{equation}
  \mathrm{halo\ overhead}\approx
  \frac{|\ghost_r|N_p}{|\owned_r|N_p}
  =
  \frac{|\ghost_r|}{|\owned_r|}.
\end{equation}
The ratio is written in \code{partition\_metadata.csv}.  The overhead is
largest for small partitions, high surface-to-volume partition shapes, and
large rank counts at fixed global problem size.

Persistent communication buffers remove per-RHS allocation for production
Maxwell halo exchange.  Surface flux kernels use reusable
\code{MaxwellSurfaceWorkspace} objects and precomputed face lifting operators.
The implementation nevertheless has some known scaling limits:
\begin{itemize}
  \item nonblocking MPI calls are currently followed by \code{MPI.Waitall}
        before computation, so there is no communication/computation overlap;
  \item root-based mesh distribution can bottleneck setup for very large
        meshes unless the distributed mesh cache or collective preparation
        path is used;
  \item diagnostics based on high-order cubature can be expensive because
        many invariant and exact-error quantities are evaluated at every
        diagnostic sample;
  \item ParaView output writes high-order Lagrange tetrahedra for owned
        elements, which is useful for visualization but can dominate storage
        for frequent output.
\end{itemize}
These are implementation tradeoffs, not correctness defects.  They also
identify the natural next optimization targets: overlap interior volume work
with halo exchange, reduce global setup bottlenecks, and tune diagnostic and
output cadence for production runs.

\section{Implemented Production Workflows}

The main distributed Maxwell frontends are:
\begin{itemize}
  \item \path{examples/distributed_poisson_bracket_maxwell.jl}: cavity
        Poisson-bracket Maxwell run with PEC or PMC boundaries, centered or
        alternating flux subject to the implemented boundary restrictions,
        optional nonlinear PML, diagnostics, checkpoint/restart, and ParaView
        time series.
  \item \path{examples/distributed_periodic_poisson_bracket_maxwell.jl}:
        periodic plane-wave run on axis-aligned translationally matching box
        meshes, with automatic or user-supplied partitions, distributed
        periodic exchange, diagnostics, checkpoint/restart, and ParaView
        output.
  \item \path{examples/distributed_poisson_bracket_pml_validation.jl}:
        production PML validation driver with mesh/flux/PML parameter sweeps,
        reflection diagnostics, metadata, and ParaView output.
  \item \path{examples/distributed_metallic_sphere_scattering.jl}: PEC sphere
        scattered-field driver with incident plane wave, nonlinear PML,
        incident-aware PEC boundary flux, PEC boundary residual diagnostics,
        optional RCS extraction, and ParaView output of scattered/incident
        and total fields.
  \item \path{examples/distributed_pec_almond_scattering.jl} and
        \path{examples/distributed_pec_almond_austin_monostatic_rcs.jl}:
        NASA PEC almond scattering/RCS workflows using the same distributed
        Maxwell, incident-boundary, PML, residual, and RCS infrastructure.
\end{itemize}

\section{What Is Not Implemented or Not Yet General}

The following limitations are visible from the inspected implementation:
\begin{itemize}
  \item dynamic load balancing is not implemented; partitions are fixed for a
        run;
  \item communication/computation overlap is not implemented in the production
        RHS path;
  \item checkpoint restart requires the same MPI rank count, DG order, and
        element partition;
  \item distributed periodic pairing requires matching opposite boundary
        topology; arbitrary nonmatching periodic meshes are not supported;
  \item Poisson-bracket alternating flux with physical PMC or absorbing
        boundaries is not implemented;
  \item the older \path{src/rhs.jl} MPI RHS prototype is not part of the
        module include chain and should not be treated as the production
        solver path.
\end{itemize}

\section{Conclusion}

DiscoGMPI implements distributed-memory three-dimensional DG Maxwell time
integration by combining a serial tetrahedral DG kernel with an owned-plus-halo
mesh representation and nearest-neighbor MPI trace exchange.  The main design
choice is pragmatic: ghost elements allow the serial face machinery to remain
unchanged, while ownership masks prevent double-counting and prevent ghost
states from advancing.  Global conservation diagnostics, convergence studies,
PML validation, and RCS workflows are all built on the same pattern: compute
rank-local owned contributions and reduce or gather them explicitly.

The current implementation is therefore already suitable for research-grade
MPI validation studies of PEC/PMC cavities, periodic plane waves, nonlinear
PML behavior, and PEC scattering/RCS workflows.  The remaining engineering
work is mostly about scalability and maintainability rather than establishing
the basic distributed numerical model: removing setup bottlenecks, overlapping
MPI exchange with computation, keeping drivers thin, and making the validation
matrix the strict release gate before large algorithmic changes.

\end{document}
